\documentclass[11pt]{article}
\usepackage[a4paper,margin=1in]{geometry}
\usepackage{amsmath}
\usepackage{booktabs}

\title{A Synthetic Example of Run Comparison}
\author{Aster Research Demo}
\date{October 2026}

\begin{document}
\maketitle

\begin{abstract}
This document is demonstration content. Its run names and measurements are synthetic;
they are not claims about a trained model. It shows how a manuscript can remain beside
run evidence while an author edits and reviews individual passages.
\end{abstract}

\section{Research question}
How should a researcher compare two runs when the observed metric and the evidence
behind it can differ? A useful workspace keeps the source manuscript, run records,
and agent conversation linked without treating an agent response as an experimental result.

\section{Method}
We illustrate a baseline and a rank-16 configuration. For a metric sequence
$m_1,\ldots,m_n$, the displayed last observation is
\begin{equation}
  m_{\mathrm{last}} = m_n.
\end{equation}
This quantity is not automatically a final evaluation score. A comparison is only
meaningful when dataset revision, evaluation split, units, and aggregation match.

\section{Illustrative results}
\begin{table}[h]
  \centering
  \begin{tabular}{lcc}
    \toprule
    Run & Configuration & Evidence status \\
    \midrule
    Baseline & reference & demo only \\
    Rank-16 & ablation & demo only \\
    \bottomrule
  \end{tabular}
  \caption{Synthetic runs used to demonstrate the writing workflow.}
\end{table}

\section{Limitations}
No statistical claim follows from a single example run. Code provenance and
evaluation metadata must be checked before reporting a difference as a result.
An author can select this paragraph, ask an agent for a clearer formulation,
review the proposed text, and then decide whether to edit the manuscript.

\end{document}
